\documentclass{article}

\usepackage{amsmath}
\usepackage{booktabs}
\usepackage{array}
\usepackage{longtable}
\usepackage{xcolor}
\usepackage{listings}
\usepackage[hidelinks]{hyperref}

\setlength{\emergencystretch}{3em}
\newenvironment{card}{\par\smallskip\noindent\begin{minipage}{\linewidth}}{\end{minipage}\par}
\lstset{basicstyle=\ttfamily\footnotesize, frame=single, keepspaces=true,
        showstringspaces=false}
\lstdefinestyle{sig}{frame=none, backgroundcolor=\color{gray!12},
        aboveskip=5pt, belowskip=2pt}

\title{UKRAINIX WATERDRONE\\[2mm]
       Technical Requirements: \texttt{core} and \texttt{parsing}\\[2mm]
       \large analysis-app, first implementation step}
\author{UKRAINIX WATERDRONE project}
\date{2 November 2026}

\begin{document}

\maketitle

\noindent This document tells the person implementing the first two packages of
\path{analysis-app/src/analysis_app/} what to write: which files and
functions, what goes in and out, and what may be assumed to exist already. It
covers \path{core/} and \path{parsing/} only; nothing else from the planned
layout is considered.

% =====================================================================
\section{Scope and ground rules}

\subsection{Files to write}

{\small
\begin{longtable}{@{}p{0.30\textwidth}p{0.64\textwidth}@{}}
\toprule
\textbf{File} & \textbf{Content}\\
\midrule
\endhead
\path{core/types.py} & Column-name constants, the \path{Table} alias, three
  data containers, four exceptions (no functions).\\
\path{core/config.py} & Path helper, text-dictionary reader, defaults access
  (5 functions).\\
\path{core/functional.py} & Two small table helpers (2 functions).\\
\path{core/__init__.py} & Re-exports every public name of the three files.\\
\path{parsing/coordinates.py} & Latitude/longitude to local metres and back,
  validity check, distance, error estimate (6 functions).\\
\path{parsing/csv_parser.py} & CSV reading, column detection, row validation,
  building the sample (9 functions).\\
\path{parsing/__init__.py} & Re-exports every public name of the two files.\\
\bottomrule
\end{longtable}
}

\subsection{Ground rules}

\begin{itemize}
\item \textbf{Functional style.} Plain functions, no classes except
      \path{NamedTuple} containers and exception classes. A function never
      modifies its inputs; it returns new objects.
\item \textbf{File access only where named.} Only \path{read_text_dictionary},
      \path{load_defaults}, \path{read_csv} and \path{load_sample} touch the
      disk. Everything else is pure.
\item \textbf{NumPy arrays are the primary data type.} All sample data are
      one-dimensional arrays of shape \path{(n,)}: \path{float64}, except
      timestamps (\path{int64}). A missing value is \path{nan}.
\item \textbf{Units.} Angles are in degrees at every public interface
      (radians only inside a function). Lengths are in metres.
\item \textbf{Type hints and docstrings} on every function: one summary line,
      then arguments, return value and raised exceptions.
\item \textbf{Errors} are raised as the project exceptions of
      \path{core/types.py}, never as \path{print} or \path{assert}.
\item \textbf{Imports.} Other code imports from the package
      (\path{analysis_app.core}, \path{analysis_app.parsing}), never from a
      file inside it. \path{core} imports nothing from \path{parsing}.
\end{itemize}

Dependencies inside this task: \path{core/functional.py} and
\path{core/config.py} use \path{core/types.py}; \path{parsing/coordinates.py}
uses \path{core}; \path{parsing/csv_parser.py} uses \path{core} and
\path{parsing/coordinates.py}. \path{core/types.py} imports nothing from the
project.

% =====================================================================
\section{What is assumed to be available}
\label{sec:assumed}

Python 3.10 or newer (assumption; the code uses \path{X | None} annotations),
NumPy, and the standard library. \textbf{Not available:} pandas, scipy,
matplotlib, and every other package of the app (normalization, isolines,
physics, inversion, pipeline, viz, gui).

The building blocks below are the only external calls the two packages need.
Their behaviour was checked with NumPy 2.5.

{\small
\begin{longtable}{@{}p{0.40\textwidth}p{0.55\textwidth}@{}}
\toprule
\textbf{Call} & \textbf{What it returns}\\
\midrule
\endhead
\path{csv.reader(f)} & Iterator over rows; each row is a list of
  \path{str}.\\
\path{np.genfromtxt(path, delimiter=",", skip_header=1, encoding="utf-8")} &
  2-D \path{float64} array (rows $\times$ columns). Empty and non-numeric cells
  become \path{nan}. One data row gives shape \path{(n_cols,)}; a header-only
  file gives shape \path{(0,)}; a row with the wrong number of columns raises
  \path{ValueError}.\\
\path{np.atleast_2d(a)} & \path{a} with at least two dimensions (turns
  \path{(n_cols,)} into \path{(1, n_cols)}).\\
\path{np.isnan(a)}, \path{np.isfinite(a)} & Boolean array, same shape.\\
\path{np.argsort(a, kind="stable")} & Integer indices that sort \path{a};
  equal values keep their order.\\
\path{np.diff(a)} & \path{a[1:] - a[:-1]}, length $n-1$.\\
\path{np.hypot(dx, dy)} & Element-wise $\sqrt{dx^2+dy^2}$.\\
\path{np.median(a)} & Median as a float.\\
\path{np.radians}, \path{np.degrees}, \path{np.sin}, \path{np.cos},
  \path{np.arcsin}, \path{np.sqrt}, \path{np.clip} & Element-wise maths on
  arrays.\\
\path{Path(__file__).resolve().parents[k]} & The $k$-th parent folder of the
  current file, as a \path{Path}.\\
\path{a.flags.writeable = False} & Makes an array read-only; assigning to an
  element then raises \path{ValueError}.\\
\path{s.isupper()} & \path{True} if all letters of \path{s} are capitals and
  there is at least one: \path{"PARAM_A"} is \path{True}, \path{"Param"} is
  \path{False}.\\
\bottomrule
\end{longtable}
}

% =====================================================================
\section{Package \texttt{core}}

\subsection{\texttt{types.py}}

No functions. Define the following.

\paragraph{Constants.}
\path{COL_TIMESTAMP = "TIMESTAMP"}, \path{COL_LATITUDE = "LATITUDE"},
\path{COL_LONGITUDE = "LONGITUDE"}, \path{COL_DEPTH = "Z"};
\path{REQUIRED_COLUMNS} = (timestamp, latitude, longitude);
\path{FIXED_COLUMNS} = \path{REQUIRED_COLUMNS} plus the depth column.

\paragraph{Alias.}
\path{Table = dict[str, np.ndarray]}: CSV column name $\to$ one-dimensional
array; all arrays have the same length; key order is the file order.

\paragraph{Containers} (\path{NamedTuple}, no methods):

{\small
\begin{longtable}{@{}p{0.17\textwidth}p{0.77\textwidth}@{}}
\toprule
\textbf{Name} & \textbf{Fields}\\
\midrule
\endhead
\path{RefOrigin} & \path{lat0_deg: float}, \path{lon0_deg: float}\\
\path{Sample} & \path{source: str}; \path{t: ndarray} (int64);
  \path{lat_deg}, \path{lon_deg}, \path{x}, \path{y}: \path{ndarray} (float64,
  \path{x}, \path{y} in metres); \path{depth: ndarray | None};
  \path{values: dict[str, ndarray]} (parameter name $\to$ readings);
  \path{ref_origin: RefOrigin}\\
\path{ParseReport} & \path{source: str}; \path{n_rows_read},
  \path{n_dropped_nan}, \path{n_dropped_bad_coordinates},
  \path{n_dropped_duplicate_time}, \path{n_dropped_glitch},
  \path{n_rows_kept}: \path{int}; \path{parameters: tuple[str, ...]};
  \path{ignored_columns: tuple[str, ...]}\\
\bottomrule
\end{longtable}
}

\noindent Invariants of a \path{Sample}: at least one point; all arrays have
the same length; \path{t} is strictly increasing; no \path{nan} anywhere;
every array is read-only.

\paragraph{Exceptions.} \path{WaterdroneError} (derived from \path{Exception}) is the base class;
\path{CsvFormatError}, \path{ConfigError} and \path{CoordinateError} derive
from it.

\subsection{\texttt{config.py}}

Assumes the repository layout \path{analysis-app/src/analysis_app/core/config.py}.

\begin{card}
\begin{lstlisting}[style=sig]
def app_root() -> Path
def config_dir() -> Path
\end{lstlisting}
\textbf{Does.} \path{app_root} returns the \path{analysis-app} folder,
computed from \path{__file__} (\path{parents[3]}), so it does not depend on the
working directory. \path{config_dir} returns \path{app_root() / "config"}.
\textbf{Returns.} absolute \path{Path}. \textbf{Raises.} nothing.
\end{card}

\begin{card}
\begin{lstlisting}[style=sig]
def read_text_dictionary(path: Path) -> dict[str, str]
\end{lstlisting}
\textbf{Does.} Reads a text file made of lines \texttt{KEY value}: the key is the
first word, the value is the rest of the line with surrounding spaces removed.
Blank lines and lines starting with \path{#} are skipped. Values stay strings.
\textbf{Returns.} dictionary in file order. \textbf{Raises.}
\path{ConfigError} if the file is missing, a line has a key but no value, or a
key appears twice.
\end{card}

\begin{card}
\begin{lstlisting}[style=sig]
def load_defaults(path: Path | None = None) -> dict[str, str]
\end{lstlisting}
\textbf{Does.} Reads the file \path{defaults.txt} in \path{config_dir()} (or the
file \path{path}) with \path{read_text_dictionary}. \textbf{Returns.} the dictionary, or an empty dictionary
if the default file does not exist (the file is optional; it is only planned).
\textbf{Raises.} \path{ConfigError} for a malformed file.
\end{card}

\begin{card}
\begin{lstlisting}[style=sig]
def get_float(settings: Mapping[str, str], key: str,
              fallback: float) -> float
\end{lstlisting}
\textbf{Does.} Converts \path{settings[key]} to a float. \textbf{Returns.}
\path{fallback} if the key is missing. \textbf{Raises.} \path{ConfigError} if
the value is not a number.
\end{card}

\subsection{\texttt{functional.py}}

\begin{card}
\begin{lstlisting}[style=sig]
def table_length(table: Table) -> int
\end{lstlisting}
\textbf{Does.} Returns the common length of all columns; \path{0} for a table
without columns. \textbf{Raises.} \path{ValueError} if the columns have
different lengths.
\end{card}

\begin{card}
\begin{lstlisting}[style=sig]
def select_rows(table: Table, selector: np.ndarray) -> Table
\end{lstlisting}
\textbf{Does.} Applies one row selector to every column: \path{selector} is a
boolean mask of length $n$ or an array of integer row indices. \textbf{Returns.}
a new table with the same keys in the same order; the input is not changed.
\textbf{Raises.} \path{ValueError} if a boolean mask has the wrong length.
\end{card}

% =====================================================================
\section{Package \texttt{parsing}}

\subsection{\texttt{coordinates.py}}
\label{sec:coord}

Constant: \path{EARTH_RADIUS_M = 6_371_000.0}. All functions accept arrays of
any length, including length 0, unless stated otherwise. Let $R$ be the Earth
radius, $\varphi,\lambda$ latitude and longitude in radians and
$(\varphi_0,\lambda_0)$ the reference origin. The projection is the simple
small-area one
\begin{equation*}
  x = R\cos\varphi_0\,(\lambda-\lambda_0), \qquad y = R\,(\varphi-\varphi_0),
\end{equation*}
and its inverse is $\varphi=\varphi_0+y/R$,
$\lambda=\lambda_0+x/(R\cos\varphi_0)$. This is a stand-in: the final formula
and error estimate follow the team's T3 lecture notes. It must therefore live
in exactly \emph{one} function pair (\path{latlon_to_local} and
\path{local_to_latlon}), so it can be replaced without touching other code.
Tracks crossing the 180$^\circ$ meridian or a pole are not supported.

\begin{card}
\begin{lstlisting}[style=sig]
def valid_latlon_mask(lat_deg: np.ndarray,
                      lon_deg: np.ndarray) -> np.ndarray
\end{lstlisting}
\textbf{Returns.} boolean array: \path{True} where both values are finite and
$-90\le\text{lat}\le90$, $-180\le\text{lon}\le180$.
\end{card}

\begin{card}
\begin{lstlisting}[style=sig]
def choose_reference_origin(lat_deg: np.ndarray,
                            lon_deg: np.ndarray) -> RefOrigin
\end{lstlisting}
\textbf{Does.} Uses the \emph{median} latitude and the median longitude (the
median is not distorted by a stray GPS point). \textbf{Raises.}
\path{CoordinateError} for empty or non-finite input.
\end{card}

\begin{card}
\begin{lstlisting}[style=sig]
def latlon_to_local(lat_deg, lon_deg, origin: RefOrigin
                    ) -> tuple[np.ndarray, np.ndarray]
\end{lstlisting}
\textbf{Returns.} \path{(x_m, y_m)}, same length as the input.
\textbf{Raises.} \path{CoordinateError} if the two inputs have different
shapes or contain non-finite values.
\end{card}

\begin{card}
\begin{lstlisting}[style=sig]
def local_to_latlon(x_m, y_m, origin: RefOrigin
                    ) -> tuple[np.ndarray, np.ndarray]
\end{lstlisting}
\textbf{Returns.} \path{(lat_deg, lon_deg)}; exact inverse of
\path{latlon_to_local}. Same \textbf{Raises}.
\end{card}

\begin{card}
\begin{lstlisting}[style=sig]
def great_circle_distance_m(lat1_deg, lon1_deg,
                            lat2_deg, lon2_deg) -> np.ndarray
\end{lstlisting}
\textbf{Does.} Haversine distance on a sphere of radius $R$, element-wise
(arguments may be arrays or scalars that broadcast):
\begin{equation*}
  a=\sin^2\tfrac{\Delta\varphi}{2}
    +\cos\varphi_1\cos\varphi_2\sin^2\tfrac{\Delta\lambda}{2},\qquad
  d = 2R\arcsin\sqrt{a}.
\end{equation*}
Clip $a$ to $[0,1]$ before the square root, because rounding can push it
slightly above 1.
\end{card}

\begin{card}
\begin{lstlisting}[style=sig]
def projection_error_estimate(lat_deg, lon_deg,
                              origin: RefOrigin) -> float
\end{lstlisting}
\textbf{Does.} Largest absolute difference, in metres, between the local-plane
distance $\sqrt{x_i^2+y_i^2}$ to the origin and the great-circle distance to
the origin, over all points. A simple stand-in for the T3 error derivation.
\textbf{Returns.} \path{0.0} for empty input.
\end{card}

\subsection{\texttt{csv\_parser.py}}

Constants: \path{DEFAULT_MAX_SPEED_MPS = 20.0} (a placeholder, see
Section~\ref{sec:open}), \path{SETTING_MAX_SPEED = "max_speed_mps"}.

\paragraph{Table convention.} Inside this file a \path{Table} holds
\path{float64} columns until \path{build_sample}. Every \path{drop_*} and
\path{sort_by_time} function takes a table and returns a new table; it must
work for tables with 0 or 1 rows.

\begin{card}
\begin{lstlisting}[style=sig]
def read_csv(path: Path) -> Table
\end{lstlisting}
\textbf{Does.} Opens the file with \path{encoding="utf-8-sig"} (removes the
byte-order mark that spreadsheet programs add) and \path{newline=""}; reads the
header with \path{csv.reader}, strips spaces from the names; reads the data
with \path{np.genfromtxt} (Section~\ref{sec:assumed}). \textbf{Returns.} table with all
columns of the file, as \path{float64}. \textbf{Raises.} \path{CsvFormatError}
for an empty header, a duplicate column name, no data rows, or a data row with
the wrong number of columns. A missing file raises the usual
\path{FileNotFoundError}.
\end{card}

\begin{card}
\begin{lstlisting}[style=sig]
def detect_columns(columns: Iterable[str]
                   ) -> tuple[tuple[str, ...], tuple[str, ...]]
\end{lstlisting}
\textbf{Does.} A \emph{parameter} is a column that is not in
\path{FIXED_COLUMNS} and whose name satisfies \path{name.isupper()}. Any other
column is \emph{ignored}. \textbf{Returns.} \path{(parameters, ignored)}, both
in file order. \textbf{Raises.} \path{CsvFormatError} if a required column is
missing (name it in the message) or if there is no parameter at all.
\end{card}

\begin{card}
\begin{lstlisting}[style=sig]
def drop_nan_rows(table: Table) -> Table
\end{lstlisting}
\textbf{Does.} Removes every row in which \emph{any} column is \path{nan}.
(Callers pass only the needed columns, so text in an ignored column cannot
remove rows.)
\end{card}

\begin{card}
\begin{lstlisting}[style=sig]
def drop_bad_coordinates(table: Table) -> Table
\end{lstlisting}
\textbf{Does.} Removes rows for which \path{valid_latlon_mask} is \path{False}.
\end{card}

\begin{card}
\begin{lstlisting}[style=sig]
def sort_by_time(table: Table) -> Table
\end{lstlisting}
\textbf{Does.} Orders all rows by \path{TIMESTAMP}, ascending, with a stable
sort.
\end{card}

\begin{card}
\begin{lstlisting}[style=sig]
def drop_duplicate_timestamps(table: Table) -> Table
\end{lstlisting}
\textbf{Does.} Keeps only the first row of each timestamp. \textbf{Requires.}
a table already sorted by time (state this in the docstring).
\end{card}

\begin{card}
\begin{lstlisting}[style=sig]
def drop_gps_glitches(table: Table, origin: RefOrigin,
                      max_speed_mps: float,
                      timestamp_unit_s: float = 1.0) -> Table
\end{lstlisting}
\textbf{Does.} Converts positions to local metres with \path{origin} and
removes isolated single-point spikes. \textbf{Requires.} sorted, unique
timestamps. With $v_k$ the speed between point $k$ and $k+1$ (distance divided
by $\Delta t\cdot$\path{timestamp_unit_s}), let $h_k = (v_k > v_{\max})$ for
$k=0,\dots,n-2$. A point $i$ is a glitch if
\begin{itemize}
\item $0<i<n-1$: $h_{i-1}$ and $h_i$ are both true;
\item $i=0$: $h_0$ is true and $h_1$ is false;
\item $i=n-1$: $h_{n-2}$ is true and $h_{n-3}$ is false.
\end{itemize}
Tables with fewer than 3 rows are returned unchanged. Known limitation: two bad
points in a row are not detected; say so in the docstring.
\end{card}

\begin{card}
\begin{lstlisting}[style=sig]
def build_sample(table: Table, parameters: Sequence[str],
                 origin: RefOrigin, source: str) -> Sample
\end{lstlisting}
\textbf{Requires.} a table that contains the required columns and all
\path{parameters}, has no \path{nan}, and is sorted with unique timestamps.
\textbf{Does.} Casts \path{TIMESTAMP} to \path{int64}; computes \path{x}, \path{y}
with \path{latlon_to_local}; takes \path{depth} from the \path{Z} column if it
exists, else \path{None}; collects \path{values} for the parameters only; sets
every array read-only. \textbf{Raises.} \path{CsvFormatError} if no rows are
left or if any \path{TIMESTAMP} is not a whole number.
\end{card}

\begin{card}
\begin{lstlisting}[style=sig]
def load_sample(path: Path, *, ref_origin: RefOrigin | None = None,
                max_speed_mps: float | None = None,
                timestamp_unit_s: float = 1.0
                ) -> tuple[Sample, ParseReport]
\end{lstlisting}
\textbf{Does.} The one function other packages call. Steps, in this order:
\begin{enumerate}
\item \path{read_csv}; \path{detect_columns}.
\item Keep only the required columns, the depth column if present, and the
      parameters.
\item \path{drop_nan_rows}, \path{drop_bad_coordinates}, \path{sort_by_time},
      \path{drop_duplicate_timestamps}.
\item If \path{ref_origin} is \path{None}, choose it with
      \path{choose_reference_origin} from the remaining rows. (A caller that
      loads a second file passes the first sample's origin so both overlay.)
\item If \path{max_speed_mps} is \path{None}, take it with \path{get_float}
      from \path{load_defaults()}, falling back to
      \path{DEFAULT_MAX_SPEED_MPS}. Then \path{drop_gps_glitches}.
\item \path{build_sample}.
\end{enumerate}
\textbf{Returns.} the sample and a \path{ParseReport}. Each \path{n_dropped_*}
is the difference of \path{table_length} before and after the corresponding
step; \path{n_rows_read} is the number of data rows in the file.
\end{card}

% =====================================================================
\section{How to check the work}

Testing for the whole app uses the repository-level \path{samples/} folder; the
exact subfolder for the hand-made files below is still to be agreed. A short
script that calls the functions and compares with the expected numbers is
enough. The work is done when:

\begin{enumerate}
\item \textbf{Round trip.} Converting a few points with \path{latlon_to_local}
      and back with \path{local_to_latlon} returns the input within
      $10^{-9}$ degrees; the origin itself maps to $(0,0)$.
\item \textbf{Clean file.} A hand-made file of about 100 clean rows loads with
      nothing dropped and \path{n_rows_kept == 100}.
\item \textbf{Defects.} A hand-made file of about 15 rows with one of each
      defect (a \path{nan} cell, a text cell, latitude 95, a repeated
      timestamp, two rows in the wrong order, one GPS spike in the middle)
      gives the expected \path{ParseReport} counts and a \path{Sample} with
      strictly increasing \path{t}.
\item \textbf{Ignored columns.} A lower-case column is not a parameter and
      appears in \path{ignored_columns}; text inside it does not drop rows.
\item \textbf{Errors.} Missing \path{LATITUDE}, header only, no parameter
      column and a duplicated column name each raise \path{CsvFormatError}.
\item \textbf{Read-only.} \path{sample.x[0] = 0.0} raises \path{ValueError};
      input tables are never modified.
\item \textbf{Config and imports.} A dictionary file with comments, blank lines
      and a duplicated key behaves as specified; a missing \path{defaults.txt}
      gives an empty dictionary; \path{Sample} is importable from
      \path{analysis_app.core} and \path{load_sample} from
      \path{analysis_app.parsing}.
\end{enumerate}

% =====================================================================
\section{Assumptions and open points}
\label{sec:open}

\begin{itemize}
\item \textbf{Timestamp unit.} The CSV says only ``integer, computer time''.
      The speed check assumes seconds; \path{timestamp_unit_s} exists so the
      unit can be changed without editing code.
\item \textbf{Maximum speed.} \path{DEFAULT_MAX_SPEED_MPS = 20.0} is a
      placeholder. Replace it with a value based on the drone's real top
      speed, or set \path{max_speed_mps} in \path{defaults.txt}.
\item \textbf{NaN policy.} One \path{nan} in any used column removes the whole
      row, including the other parameters of that row. Revisit if real data
      has many isolated gaps.
\item \textbf{Text file format.} \texttt{KEY value} lines are a proposal; the
      physics dictionary will reuse the reader, so confirm it before more
      files depend on it.
\item \textbf{Projection.} The simple formula is provisional (see
      Section~\ref{sec:coord}).
\end{itemize}

\subsection*{Deliberately not included}

Planned items that these two packages do not need are \emph{not} part of this
task and are added when the module that needs them is written:
\path{compose}/\path{pipe} and per-parameter mapping in \path{functional.py};
\path{samples_dir()} in \path{config.py}; the containers for mesh, isoline,
origin, model spec and normalization result in \path{types.py}.

% =====================================================================
\section{Where to read (for beginners)}
{\small

\subsection*{NumPy}
\begin{itemize}
\item NumPy: the absolute basics for beginners (arrays, shape, dtype,
      indexing, sorting):\\
      \url{https://numpy.org/doc/stable/user/absolute_beginners.html}
\item Indexing on ndarrays; read \emph{Boolean array indexing}, the basis of
      \path{select_rows} and all \path{drop_*} functions:\\
      \url{https://numpy.org/doc/stable/user/basics.indexing.html}
\item \path{numpy.genfromtxt}, for \path{read_csv}:\\
      \url{https://numpy.org/doc/stable/reference/generated/numpy.genfromtxt.html}
\item Real Python, NumPy Tutorial: Your First Steps Into Data Science in Python
      (vectorization, filtering, sorting):\\
      \url{https://realpython.com/numpy-tutorial/}
\item For \path{argsort}, \path{diff}, \path{hypot}, \path{median},
      \path{isnan} and \path{clip}, use the search box of
      \url{https://numpy.org/doc/stable/} and read the function's page.
\end{itemize}

\subsection*{Python}
\begin{itemize}
\item \path{csv} module (\path{csv.reader}, why \path{newline=""}):\\
      \url{https://docs.python.org/3/library/csv.html}
\item \path{pathlib} (\path{Path}, \path{resolve}, \path{parents},
      \path{exists}):\\
      \url{https://docs.python.org/3/library/pathlib.html}
\item \path{typing} module; look for \path{NamedTuple} on the page:\\
      \url{https://docs.python.org/3/library/typing.html}
\item Errors and Exceptions (\path{try}/\path{except}, \path{raise}, your own
      exception classes):\\
      \url{https://docs.python.org/3/tutorial/errors.html}
\end{itemize}

\subsection*{Coordinates and distances}
\begin{itemize}
\item Movable Type Scripts, distance between two latitude/longitude points
      (haversine formula, equirectangular approximation, with code):\\
      \url{https://www.movable-type.co.uk/scripts/latlong.html}
\item Haversine formula on Wikipedia (derivation, clipping note):\\
      \url{https://en.wikipedia.org/wiki/Haversine_formula}
\item The team's own T3 lecture notes (kept outside the repository):
      \emph{coordinate system lecture}, \emph{geographic coordinates
      translation} and \emph{error value derivation}.
\end{itemize}

}
\end{document}
